\documentclass[11pt]{article}
\usepackage[margin=1in]{geometry}
\usepackage{amsmath,amssymb}
\usepackage{graphicx}
\usepackage{booktabs}
\usepackage{hyperref}
\usepackage{xcolor}
\usepackage{caption}
\usepackage{subcaption}
\usepackage{multirow}

\title{When Do Transformers Learn Structure?\\Affine Maps and Compositional Generalization\\in Synthetic Relational Tasks}
\date{}

\begin{document}
\maketitle

\begin{abstract}
Do transformers learn structured representations of relations, or do they memorize
input--output associations? We investigate this question through two complementary
synthetic studies. In the first (Salzburg), we test whether a single-layer transformer
implements relations as affine maps between subject and object embeddings. In the
second (Graz), we test whether a transformer can generalize relational predictions to
unseen entity pairings by learning compositional attribute structure. Both studies
reveal the same finding: \emph{memorization is the default}. Structured solutions ---
affine relation maps, compositional type encodings --- emerge only when capacity is
constrained to a problem-dependent sweet spot. Too little capacity and no good
structured solution exists; too much and memorization is cheaper than structure. In
Salzburg, a tiny auxiliary loss (0.02\% of the main objective) suffices to nudge the
model from 97\% to 100\% affine accuracy, showing that structured solutions are
reachable but not found by standard training. In Graz, the optimal hidden dimension
for generalization follows the same non-monotonic pattern: dim=16 achieves 90\%
held-out accuracy while dim=64 achieves only 67.5\%. These results suggest that the
memorization--structure tradeoff is governed by the ratio of model capacity to task
complexity, with implications for understanding when large language models learn
generalizable representations.
\end{abstract}

\section{Introduction}

Large language models encode vast amounts of relational knowledge --- facts about
entities, their properties, and their relationships. But the nature of this encoding
remains debated: do models learn \emph{structured} representations that support
generalization, or do they \emph{memorize} specific associations?

Recent work has provided evidence for linear structure in LLM representations.
Hernandez et al.~\cite{hernandez2024} showed that certain relations are approximately
implemented as linear maps between subject and object embeddings. However, it remains
unclear when such structure emerges and when models instead rely on memorization.

We approach this question through two synthetic studies designed for precise
experimental control:

\begin{itemize}
\item \textbf{Salzburg}: Can a transformer learn to implement relations as affine maps
  from subject to object embeddings? We vary the hidden dimension and measure whether
  the internal computation is well-approximated by per-relation affine transformations.
\item \textbf{Graz}: Can a transformer generalize relational predictions to unseen
  entity pairings by learning compositional attribute structure? Entities are defined by
  combinatorial attribute types, and relations select on shared attributes. We test
  whether the model discovers this compositional rule or memorizes specific pairings.
\end{itemize}

Both studies converge on the same finding: structured solutions emerge at an
intermediate capacity --- a sweet spot where memorization is too expensive but
sufficient representational capacity exists for the structured solution. This
non-monotonic relationship between capacity and structure is the central contribution
of this paper.

\section{Related Work}

\paragraph{Linear relational structure.}
Hernandez et al.~\cite{hernandez2024} demonstrated that relations in large language
models can be approximated by linear maps between subject and object representations,
proposing the delta-cosine similarity metric ($\Delta\cos$) to quantify this structure.
Our Salzburg experiments extend this analysis to controlled settings where we can
manipulate capacity and measure the conditions under which affine structure emerges.

\paragraph{Grokking and delayed generalization.}
Power et al.~\cite{power2022} showed that neural networks can transition from
memorization to generalization long after fitting the training data, a phenomenon termed
``grokking.'' Our results are complementary: we study the \emph{conditions} under which
structure emerges, rather than its temporal dynamics.

\paragraph{Compositional generalization.}
The ability to generalize compositionally --- combining known primitives in novel ways
--- has been studied extensively in sequence-to-sequence models~\cite{lake2018,kim2020}.
Our Graz experiments test a specific form of compositionality: whether attribute-level
structure supports generalization to unseen entity-level pairings.

\paragraph{Information bottleneck.}
The information bottleneck principle~\cite{tishby2000} suggests that compressed
representations retain task-relevant information while discarding noise. Our dimension
sweep results are consistent with this view: moderate compression forces structured
representations, while excess capacity permits unstructured memorization.

\section{Experimental Setup}

\subsection{Salzburg: Affine Relation Maps}

\paragraph{Data.}
The vocabulary consists of 100 subjects ($a_0, \ldots, a_{99}$), 100 objects
($b_0, \ldots, b_{99}$), 2 relation tokens ($r_0, r_1$), and a separator token ($x$).
Training sequences are 4-token triples $[x, r_k, a_i, b_j]$ where $b_j$ is the correct
object for subject $a_i$ under relation $r_k$. Relation $r_0$ is the identity mapping
($a_i \mapsto b_i$) and $r_1$ is a cyclic shift by 50 ($a_i \mapsto b_{(i+50) \bmod 100}$).

\paragraph{Model.}
A single-layer transformer with 4 attention heads, trained with next-token prediction
(cross-entropy loss). The hidden dimension $d$ varies across experiments
($d \in \{8, 16, 32, 64\}$).

\paragraph{Evaluation.}
We measure two quantities: (1)~\emph{object prediction accuracy} --- does the model
predict the correct object? (2)~\emph{affine accuracy} --- we fit a per-relation affine
map $\hat{o} = W_k \cdot h_0(a_i) + b_k$ from input embeddings to output representations
via least-squares regression on all 100 subjects, then check whether the affine-predicted
output $\hat{o}$ falls in the correct Voronoi cell (closest to the true object embedding).

\subsection{Graz: Compositional Attribute Relations}

\paragraph{Entity structure.}
Each entity is characterized by two attribute groups: an \emph{f-type} and a
\emph{g-type}. With $n=5$ features and $k=2$ attributes per group, each type is a
2-element subset of $\{0,\ldots,4\}$, giving $\binom{5}{2} = 10$ possible types. We
create $10^2 = 100$ entities, one for each (f-type, g-type) combination. Objects' types
are shifted relative to subjects' via shifts coprime to 10 ($r_0$-shift$\,=3$,
$r_1$-shift$\,=7$), ensuring distinct type assignments.

\paragraph{Relations.}
Relation $R_0$ connects subject $a_i$ to all objects $b_j$ sharing the same f-type.
Relation $R_1$ connects subject $a_i$ to all objects $b_j$ sharing the same g-type.
Each subject has 10 valid targets per relation (chance level: $10/100 = 10\%$).

\paragraph{Training data.}
Sequences of length 48 contain a mix of four sentence types:
\begin{itemize}
\item \emph{Attribute sentences} $[x, e, a]$: entity $e$ with one of its attributes $a$
  (probability ${\approx}0.27$)
\item \emph{Relation sentences} $[x, r_k, a_i, b_j]$: subject $a_i$ mapped to a valid
  object $b_j$ via relation $r_k$ (probability $0.264$)
\item \emph{Reverse lookups} $[x, M, a, e]$: marker $M \in \{S, O\}$ with attribute $a$
  mapping to entity $e$ (probability $0.264$)
\item \emph{Contrastive sentences} $[x, Q_k, b_j, a_i, y/n]$: query whether object $b_j$
  and subject $a_i$ are related via $r_k$, answered with $y$ (yes) or $n$ (no)
  (probability $0.20$)
\end{itemize}
Contrastive sentences use dedicated query markers ($Q_0, Q_1$) distinct from the
relation tokens, and reverse the entity order (object before subject) to prevent
interference with generative relation sentences.

We hold out 20\% of (subject, relation) triples at random and train on 62{,}500
sequences for 800 epochs. The model is a single-layer transformer with 4 attention heads.
The hidden dimension $d$ varies across experiments ($d \in \{8, 16, 32, 64\}$).

\paragraph{Evaluation.}
(1)~\emph{Held-out relation accuracy}: given $[x, r_k, a_i]$, does the model predict
any valid object $b_j$? Evaluated on all held-out triples. (2)~\emph{Contrastive
accuracy}: does the model correctly predict $y$ or $n$? (3)~\emph{Affine probe}: a
post-hoc linear classifier trained on half the subjects' hidden representations to
predict type membership, tested on the other half.

\section{Salzburg Results}

\subsection{Affine Structure Emerges at the Right Scale}

All models achieve near-perfect object prediction accuracy ($\geq 99.5\%$) regardless
of hidden dimension. However, \emph{affine accuracy} --- the degree to which the
internal computation is well-approximated by per-relation affine maps --- varies
dramatically (Table~\ref{tab:salzburg}).

\begin{table}[h]
\centering
\caption{Salzburg dimension sweep. All models predict correctly, but affine structure
is non-monotonic in capacity. The existence test fits affine maps to fixed random
unit-vector object embeddings via 20K Adam steps.}
\label{tab:salzburg}
\begin{tabular}{lccccl}
\toprule
& dim=8 & dim=16 & dim=32 & dim=64 & \\
\midrule
Parameters & 5{,}128 & 11{,}792 & 29{,}728 & 84{,}032 & \\
Object accuracy & 99.5\% & 100\% & 100\% & 100\% & \\
Affine accuracy (L1 avg) & 73.0\% & 97.0\% & 100\% & 27.5\% & \\
Existence test (Voronoi) & 70.0\% & 99.0\% & 100\% & --- & \\
\midrule
Regime & \multicolumn{1}{c}{(A)} & \multicolumn{2}{c}{Sweet spot} & \multicolumn{1}{c}{(C)} & \\
\bottomrule
\end{tabular}
\end{table}

Three regimes are evident:

\begin{itemize}
\item \textbf{dim=8 --- No good solution exists (A).} The existence test, which fits
  affine maps under ideal conditions (fixed unit-vector objects, 20K optimization steps),
  achieves only 70\% Voronoi accuracy. The transformer matches this at 73\%. The
  limitation is geometric: 100 objects in 8 dimensions cannot be separated well enough
  for affine maps to work.
\item \textbf{dim=16--32 --- Sweet spot.} The existence test achieves 99--100\%, and the
  transformer tracks this closely (97--100\%). Structure is both representable and found
  by SGD, because the affine map is the most parameter-efficient solution.
\item \textbf{dim=64 --- Memorization wins (C).} Despite 84K parameters for only 200
  facts (420 parameters per fact), affine accuracy collapses to 27.5\%. The model can
  dedicate unique directions to each subject--object pair without needing affine
  regularity.
\end{itemize}

Figure~\ref{fig:dim_sweeps} (left) shows this non-monotonic relationship.

\subsection{Auxiliary Loss Nudges Toward Structure}

At dim=16, the transformer achieves 97\% affine accuracy --- close to but not matching
the existence test's 99\%. To test whether the gap is due to architectural limitation or
optimization, we add a tiny auxiliary loss encouraging affine structure:
\[
\mathcal{L} = \mathcal{L}_{\text{CE}} + \lambda \sum_{k} \frac{1}{N} \sum_i \| W_k \cdot h_0(a_i) + b_k - h_1(a_i) \|^2
\]
where $W_k, b_k$ are learnable affine parameters (one per relation) optimized jointly
with the transformer. With $\lambda = 1.0$, the auxiliary loss stabilizes at
${\sim}0.0003$ --- roughly 0.02\% of the cross-entropy loss of ${\sim}1.32$.

This negligible gradient signal pushes affine accuracy from 97\% to \textbf{100\%},
exceeding even the existence test's 99\% (because the transformer jointly optimizes all
embeddings rather than using fixed random objects). The architecture can clearly
implement the affine solution; standard SGD simply does not find it without explicit
pressure.

\subsection{Mechanistic Analysis}

A logit lens analysis reveals how the transformer implements relation-dependent mappings
(Table~\ref{tab:logit_lens}).

\begin{table}[h]
\centering
\caption{Logit lens: object prediction accuracy at intermediate computation stages
(dim=16, evaluated on 32{,}000 positions).}
\label{tab:logit_lens}
\begin{tabular}{lcc}
\toprule
Stage & Description & Object accuracy \\
\midrule
$h_0$ & Embedding only & 17.2\% \\
$h_0 + \text{attn}$ & Post-attention & 14.4\% \\
$h_0 + \text{attn} + \text{MLP}$ & Full residual & 100\% \\
\bottomrule
\end{tabular}
\end{table}

Post-attention accuracy (14.4\%) is no better than the raw embedding (17.2\%), despite
the attention output being near-perfectly affine-predictable from the input
($\cos(\hat{o}, o) = 0.987$). Attention routes information --- copying subject and
relation embeddings into the prediction position --- but cannot perform the
\emph{relation-dependent} transformation needed to select the correct object. This
requires a multiplicative interaction between relation identity and subject embedding,
which only the MLP's GELU nonlinearity can provide.

The sublayer decomposition confirms this: the attention output is 0.987-affine in the
input (mechanically, attention computes weighted sums of linearly-transformed inputs),
while the MLP output is 0.773-affine. The MLP implements the essential nonlinear gating
that selects which affine map to apply based on the relation token.

\section{Graz Results}

\subsection{Capacity Controls Structure vs.\ Memorization}

We train the Graz model at four hidden dimensions, holding all other parameters fixed.
Table~\ref{tab:graz_dim} and Figure~\ref{fig:dim_sweeps} (right) show the results.

\begin{table}[h]
\centering
\caption{Graz dimension sweep (5 features, 100 entities, 10 types). All models achieve
100\% train accuracy. Chance level is 10\%.}
\label{tab:graz_dim}
\begin{tabular}{lcccc}
\toprule
& dim=8 & dim=16 & dim=32 & dim=64 \\
\midrule
Held-out accuracy & 42.5\% & \textbf{90.0\%} & 85.0\% & 67.5\% \\
$\times$ chance & 4.25$\times$ & 9.0$\times$ & 8.5$\times$ & 6.75$\times$ \\
Contrastive ($+$/$-$) & 85.2/94.0 & 88.8/97.8 & 88.8/99.2 & 90.0/97.6 \\
Affine probe (test) & 26.0\% & 50.0\% & 55.0\% & 63.0\% \\
R0 partner similarity & +0.22 & +0.43 & +0.18 & +0.10 \\
R1 partner similarity & +0.20 & +0.42 & +0.18 & +0.13 \\
\bottomrule
\end{tabular}
\end{table}

The same non-monotonic pattern as Salzburg emerges:

\begin{itemize}
\item \textbf{dim=8 --- Capacity collapse.} Embedding overlap reaches 0.91--0.93,
  meaning entities are nearly indistinguishable. The model cannot represent 100 distinct
  entities in 8 dimensions while also encoding type structure. Held-out accuracy (42.5\%)
  is well above chance but far below the structured solution.
\item \textbf{dim=16 --- Sweet spot.} Held-out accuracy peaks at \textbf{90.0\%}
  (9$\times$ chance). Partner similarity is highest (+0.43), indicating that
  relation-relevant entities are tightly clustered. The information bottleneck forces
  compact, compositional representations.
\item \textbf{dim=32--64 --- Diminishing structure.} As capacity grows, held-out
  accuracy monotonically decreases (85.0\% $\to$ 67.5\%). Partner similarity drops
  (+0.18 $\to$ +0.10): with more dimensions, the model spreads representations apart
  rather than clustering by type. Memorization becomes progressively cheaper.
\end{itemize}

\subsection{The Affine Probe Paradox}

An intriguing pattern emerges in the affine probe results: \emph{the best-generalizing
model has the worst linear readout of type information}. At dim=16, held-out accuracy is
90\% but the affine probe achieves only 50\%. At dim=64, held-out accuracy drops to
67.5\% but the affine probe improves to 63\%.

This is not a contradiction. The affine probe measures \emph{linear accessibility} of
type information, not its functional utility. At dim=16, the model encodes type
information in a compact, possibly nonlinear representation that supports accurate
relation predictions but is hard for a linear classifier to decode. At dim=64, the
model has room to encode type information in linearly separable directions --- but
also has room to \emph{not use} that information for predictions, relying instead on
memorized associations.

The same phenomenon appears in Salzburg: at dim=64, the model achieves 100\% object
accuracy but only 27.5\% affine accuracy. High-dimensional representations can encode
structure that probes detect while the model's actual computation follows a different
path.

\subsection{Embedding Structure Reflects the Bottleneck}

Figure~\ref{fig:graz_detail} (right) shows partner similarity across dimensions. At
dim=16, entities sharing a relation-relevant type have cosine similarity +0.43 ---
they are tightly clustered. At dim=64, this drops to +0.10. The tighter clustering at
dim=16 reflects the information bottleneck: the model must reuse representational
dimensions across entities, and organizing by type is the most efficient strategy.

At dim=8, partner similarity is moderate (+0.21) but overall embedding overlap is
extreme (0.91--0.93): all entities are squeezed into a tiny region of the 8-dimensional
space. The model cannot distinguish entities well enough for accurate relational
predictions.

\begin{figure}[t]
\centering
\includegraphics[width=\textwidth]{figures/fig1_dimension_sweeps.pdf}
\caption{Non-monotonic relationship between hidden dimension and structured behavior
in both Salzburg (left: affine accuracy) and Graz (right: held-out generalization).
Both show a sweet spot at intermediate dimensions, with collapse at low dimensions and
memorization at high dimensions.}
\label{fig:dim_sweeps}
\end{figure}

\begin{figure}[t]
\centering
\includegraphics[width=\textwidth]{figures/fig3_graz_dim_sweep.pdf}
\caption{Graz dimension sweep detail. (a) Held-out generalization peaks at dim=16
while affine probe accuracy increases monotonically with dimension. (b) Embedding
partner similarity --- cosine similarity between entities sharing a relation-relevant
type --- is highest at dim=16, reflecting the information bottleneck.}
\label{fig:graz_detail}
\end{figure}

\section{Discussion}

\subsection{The Memorization--Structure Tradeoff}

Both studies reveal the same governing principle: \emph{structured solutions emerge when
they are the most parameter-efficient path to low loss}. At the sweet spot (dim=16--32
for Salzburg, dim=16 for Graz), the model has enough capacity to represent the
structured solution but not enough to memorize all facts individually. Structure is not
an intrinsic property of the architecture --- it is an emergent consequence of the
capacity--complexity ratio.

This explains why the two studies show parallel results despite testing different forms
of structure (affine maps vs.\ compositional types). The mechanism is the same:
dimensional compression creates a bottleneck that makes structured representations
cheaper than brute-force memorization.

\subsection{Implications for Large Language Models}

Real LLMs are massively overparameterized relative to any individual relation. A model
with billions of parameters storing thousands of facts about a relation operates deep
in the ``memorization is cheap'' regime (dim=64 in our experiments). Our results predict
that such models would default to memorization rather than structured relational
representations.

However, two factors may push toward structure in practice:
\begin{enumerate}
\item \emph{Frequency}: Relations that appear in millions of training examples create
  stronger gradient pressure toward efficient encoding than those appearing in hundreds.
\item \emph{Implicit contrastive signals}: Natural language contains negation, contrast,
  and questions (``Is Paris the capital of Germany? No.'') that provide the kind of
  positive/negative signal our contrastive sentences supply explicitly.
\end{enumerate}

The auxiliary loss result in Salzburg is particularly suggestive: a gradient signal
comprising just 0.02\% of the total loss suffices to reorganize internal representations
from 97\% to 100\% affine. This implies that structured solutions, even in
overparameterized models, may be surprisingly close to the learned solution --- reachable
with minimal additional pressure.

\subsection{Limitations}

Our experiments use single-layer transformers with at most 100--200 entities.
Deeper models, larger vocabularies, and natural-language data may exhibit different
dynamics. The non-monotonic capacity--structure relationship may shift with scale, and
multi-layer models could develop structure at higher dimensions than our single-layer
models. The Graz contrastive results (Appendix~\ref{app:contrastive}) suggest that
training objectives also matter --- an aspect we have not fully disentangled from
capacity effects in the main experiments.

\section{Conclusion}

We have presented two complementary studies showing that transformers default to
memorization when capacity is abundant. Structured solutions --- affine relation maps
(Salzburg) and compositional type encodings (Graz) --- emerge only at a
capacity-dependent sweet spot. Too little capacity and no good structured solution
exists; too much and memorization is cheaper. This non-monotonic relationship is
consistent across both relational tasks and suggests a general principle:
\emph{the ratio of model capacity to task complexity governs whether transformers learn
structure or memorize}.

Two levers can push models toward structure: \emph{capacity constraints} (dimensional
bottlenecks) and \emph{objective-level pressure} (auxiliary losses or contrastive
training). Both are effective because they change the relative cost of memorization
versus structure, tipping the balance toward generalizable representations.

\appendix

\section{Contrastive Learning Ablation}
\label{app:contrastive}

The Graz experiments in the main body include contrastive sentences in the training
data. Here we present evidence from a larger-scale setup (7 features, 441 entities,
21 types) showing that contrastive learning is a key ingredient for generalization and
that it requires sufficient data to be effective.

\subsection{Contrastive vs.\ No Contrastive}

Table~\ref{tab:contrastive_ablation} compares the baseline (no contrastive sentences)
with contrastive training at different data scales.

\begin{table}[h]
\centering
\caption{Effect of contrastive learning (7 features, 441 entities). Contrastive
sentences hurt with insufficient data but enable strong generalization with 10$\times$
more data. Chance level is 4.8\%.}
\label{tab:contrastive_ablation}
\begin{tabular}{lccc}
\toprule
& No contrastive & Contrastive, 6.25K seq & Contrastive, 62.5K seq \\
& (phase3j, dim=32) & (dim=64) & (phase3l, dim=64) \\
\midrule
Held-out accuracy & 50.6\% & 30.7\% & \textbf{69.3\%} \\
$\times$ chance & 10.6$\times$ & 6.4$\times$ & 14.6$\times$ \\
Contrastive $+$/$-$ & --- & 59\%/53\% & 91.6\%/98.6\% \\
Affine probe (test) & 51.8\% & --- & 83.0\% \\
\bottomrule
\end{tabular}
\end{table}

With insufficient data (6{,}250 sequences), the contrastive task is barely learned
(59\%/53\%, near chance) and generalization \emph{drops} from 50.6\% to 30.7\%.
The extra sentence type competes for model capacity without providing useful signal.
With 10$\times$ more data (62{,}500 sequences), contrastive accuracy reaches
91.6\%/98.6\% and generalization jumps to 69.3\%. The affine probe improves from
51.8\% to 83.0\%, indicating that type information becomes linearly accessible.

The contrastive space is quadratic: $441 \times 441 \times 2 \approx 390$K possible
(subject, object, relation) triples. The 6{,}250-sequence dataset covers roughly 3\%
of this space; the 62{,}500-sequence dataset covers roughly 30\%.

\begin{figure}[h]
\centering
\includegraphics[width=0.55\textwidth]{figures/fig5_contrastive_scaling.pdf}
\caption{Contrastive learning at 7 features / 441 entities. Insufficient contrastive
data hurts generalization; 10$\times$ more data enables the contrastive task to drive
compositional structure.}
\label{fig:contrastive_scaling}
\end{figure}

\subsection{Exhaustive vs.\ Random Contrastive Sampling}

We also tested exhaustive enumeration: pre-generating all positive contrastive examples
(every valid (subject, relation, object) triple) with matched random negatives, then
drawing from this pool during training (100K sequences). This \emph{hurt} performance:
held-out accuracy dropped from 69.3\% (random sampling) to 56.3\% (exhaustive pool).

The likely cause is that random sampling provides a richer learning signal: different
negatives appear for each positive across epochs, creating natural variation that the
model cannot exploit through simple pattern matching. The exhaustive pool pairs each
positive with a fixed negative, reducing diversity.

\subsection{More Relation Data Does Not Help}

Increasing the fraction of relation sentences from $p = 0.33$ to $p = 0.50$ (with 10K
sequences, no contrastive) also \emph{hurt} generalization: held-out accuracy dropped
from 50.6\% to 42.6\%. More positive-only relation data enables more aggressive
memorization without providing the negative signal needed to learn the type rule.

\section{Phase Progression and Feature Reduction}
\label{app:progression}

\subsection{From Memorization to Generalization}

Table~\ref{tab:progression} shows the progression of Graz experiments from pure
memorization (Phase 2) to strong generalization (Phase 3o3).

\begin{table}[h]
\centering
\caption{Graz experimental progression. Each row represents the best result at that
configuration.}
\label{tab:progression}
\begin{tabular}{llccccc}
\toprule
Phase & Key change & Features & Entities & dim & Held-out & $\times$ chance \\
\midrule
Phase 2 & Train/test subject split & 10 & 120 & 32 & 0\% & --- \\
Phase 3 & Held-out triples & 10 & 120 & 32 & 35\% & 42$\times$ \\
Phase 3j & 7 features, $k{=}2$ & 7 & 441 & 32 & 50.6\% & 10.6$\times$ \\
Phase 3l & + contrastive, dim=64 & 7 & 441 & 64 & 69.3\% & 14.6$\times$ \\
Phase 3n & 6 features & 6 & 225 & 64 & 87.8\% & 13.2$\times$ \\
Phase 3o3 & 5 features, dim=16 & 5 & 100 & 16 & \textbf{90.0\%} & 9.0$\times$ \\
\bottomrule
\end{tabular}
\end{table}

\begin{figure}[h]
\centering
\includegraphics[width=0.85\textwidth]{figures/fig4_graz_progression.pdf}
\caption{Progression from memorization to generalization across Graz experiments.
Red: memorization-dominated. Orange: partial generalization. Green: strong
generalization driven by contrastive learning and/or capacity matching.}
\label{fig:progression}
\end{figure}

\subsection{Phase 2: Pure Memorization}

In Phase 2, 80 subjects appear in relation sentences during training and 40 are held
out entirely. The model achieves 100\% accuracy on train subjects and \emph{0\%} on
test subjects. The affine probe also shows 0\% test accuracy
($\cos(\hat{o}, o) \approx 0.04$ on test vs.\ 0.90 on train), confirming that
memorization is baked into the embedding space itself.

With 33{,}568 parameters and only 160 relation facts (80 subjects $\times$ 2 relations),
the model has ${\sim}210$ parameters per fact. Memorization is trivially cheap, and the
model has no gradient pressure to discover the compositional rule.

\subsection{Phase 3: Matrix Completion}

When all subjects appear in relations but 20\% of triples are held out, generalization
rises from 0\% to 35\%. The affine probe shows a parallel improvement (0\% $\to$
35\% test accuracy). However, the model's \emph{actual} predictions on held-out triples
remain at 0\% --- the representations carry partial relational structure that the
model's prediction head ignores in favor of the memorization shortcut.

\subsection{Feature Reduction}

Reducing the number of features from 10 to 5 progressively simplifies the type space
(from 120 types to 10), making the compositional rule more discoverable. Combined with
contrastive learning and dimension tuning, this path culminates in the 90\% held-out
accuracy reported in the main body.

\section{Salzburg: Additional Mechanistic Details}
\label{app:salzburg_mech}

\subsection{Sublayer Decomposition}

The transformer residual decomposes as $h_1 = h_0 + \text{attn\_out} + \text{mlp\_out}$.
Table~\ref{tab:sublayer} shows the affine predictability of each component from the
input embedding.

\begin{table}[h]
\centering
\caption{Sublayer decomposition (dim=16, natural context). The affine cosine similarity
$\cos(\hat{o}, o)$ measures how well each component can be predicted as an affine
function of the input embedding $h_0$.}
\label{tab:sublayer}
\begin{tabular}{lcc}
\toprule
Target & $\cos(\hat{o}, o)$ & Interpretation \\
\midrule
Attention output & 0.987 & Near-perfectly affine (mechanically linear) \\
MLP output & 0.773 & Substantially but imperfectly affine \\
Full $h_1$ residual & 0.735 & Harder to predict (sum of all components) \\
\bottomrule
\end{tabular}
\end{table}

Attention is mechanically linear: for fixed attention weights, it computes a weighted
sum of value-projected inputs, which is affine in the input. The MLP's GELU
nonlinearity introduces the relation-dependent gating needed to select different
transformations for different relations.

\subsection{Why the MLP Is Essential}

Implementing relation-dependent affine maps $a_i \mapsto W_k \cdot a_i$ requires a
multiplicative interaction between relation identity and subject embedding. Attention
combines its inputs \emph{additively} (weighted sum of value vectors) and applies the
\emph{same} value projection $W_V$ regardless of the relation. Only the MLP's
pointwise nonlinearity can implement the gating operation ``if relation = $k$, apply
$W_k$.'' This explains why post-attention accuracy (14.4\%) is no better than chance
despite the attention output being highly structured.

\begin{thebibliography}{9}

\bibitem{hernandez2024}
E.~Hernandez, A.~S.~Sharma, T.~Haklay, K.~Meng, M.~Wattenberg, J.~Andreas,
Y.~Belinkov, and D.~Bau.
\newblock Linearity of relation decoding in transformer language models.
\newblock In \emph{ICLR}, 2024.

\bibitem{power2022}
A.~Power, Y.~Burda, H.~Edwards, I.~Babuschkin, and V.~Misra.
\newblock Grokking: Generalization beyond overfitting on small algorithmic datasets.
\newblock \emph{arXiv preprint arXiv:2201.02177}, 2022.

\bibitem{lake2018}
B.~M.~Lake and M.~Baroni.
\newblock Generalization without systematicity: On the compositional skills of
  sequence-to-sequence recurrent networks.
\newblock In \emph{ICML}, 2018.

\bibitem{kim2020}
N.~Kim and A.~Linzen.
\newblock COGS: A compositional generalization challenge based on semantic
  interpretation.
\newblock In \emph{EMNLP}, 2020.

\bibitem{tishby2000}
N.~Tishby, F.~C.~Pereira, and W.~Bialek.
\newblock The information bottleneck method.
\newblock \emph{arXiv preprint physics/0004057}, 2000.

\end{thebibliography}

\end{document}
